\section{特征集定义}

\subsection{\config{manual_basic}}

\config{manual_basic} 是当前主线特征集。它把每个振动快照转换为可解释的统计特征，并分别在 horizontal、vertical 和 magnitude 通道上计算。它的设计目的不是穷尽所有信号处理方法，而是在两个数据集、三个任务之间提供低开销、可复现、可解释的共同基线。

时域特征族包括：

\begin{itemize}[leftmargin=2em]
  \item mean
  \item standard deviation
  \item root mean square
  \item mean absolute value
  \item peak-to-peak
  \item skewness
  \item kurtosis
  \item crest factor
  \item impulse factor
  \item clearance factor
\end{itemize}

频域特征族包括：

\begin{itemize}[leftmargin=2em]
  \item spectral centroid
  \item spectral bandwidth
  \item RMS frequency
  \item peak frequency
  \item spectral entropy
\end{itemize}

这些特征和任务之间的关系不是事后随意解释，而是来自轴承退化的基本物理直觉：磨损、剥落或润滑恶化往往会提高振动能量、扩大幅值分布、引入冲击成分，并改变频谱复杂度。表 \ref{tab:feature_family_interpretation} 将这一直觉和本轮实际发现放在一起。

\begin{table}[htbp]
\centering
\caption{特征族、物理含义与本轮发现}
\label{tab:feature_family_interpretation}
\scriptsize
\begin{tabularx}{\textwidth}{lYYYY}
\toprule
特征族 & 示例 & 物理含义 & 预期任务相关性 & 本轮实际发现与 caveat \\
\midrule
amplitude / energy & RMS, mean absolute, mean & 振动能量和平均幅值 & RUL, Health State, Early Fault & 两个数据集最稳定 family；\feature{mag__time__rms} 是 label-source。\\
dispersion & std, variance & 振动分散程度 & RUL 与状态分离 & \feature{mag__time__std}, \feature{h__time__std}, \feature{v__time__std} 多次进入 A/B 级。\\
peak-to-peak / shock & ptp, crest factor, impulse factor & 冲击和极值变化 & Early Fault 候选 & XJTU C2 中更明显，但未成为全局主线。\\
higher-order statistics & skewness, kurtosis & 分布形状、非高斯冲击 & Early Fault 辅助 & 本轮没有形成跨数据集稳定主结论。\\
spectral location & centroid, RMS frequency, peak frequency & 频谱能量中心与主峰位置 & 工况/故障辅助诊断 & 对转速和工况敏感；PHM Early Fault 中为 secondary。\\
spectral complexity & entropy & 频谱复杂度 & 条件性 Early Fault & XJTU C1 和 cross-condition 中出现，但属于 C 级诊断辅助。\\
tsfresh minimal & standard deviation, RMS, variance, maximum & 自动统计摘要 & 可能补充手工统计 & PHM 可运行但主要重复 amplitude/statistics；XJTU full-size blocked。\\
\bottomrule
\end{tabularx}
\end{table}


\subsection{\config{manual_tsfresh_basic}}

\config{manual_tsfresh_basic} 由手工统计特征和 tsfresh minimal 特征组成。它的作用是验证自动化统计库是否能补充 \config{manual_basic}。

PHM2012 上该配置可以完整运行，ranking 中有 20 个 \feature{tsfresh__} 特征。进入 top 区域的代表项包括：

\begin{itemize}[leftmargin=2em]
  \item \feature{tsfresh__h__standard_deviation}
  \item \feature{tsfresh__h__root_mean_square}
  \item \feature{tsfresh__h__variance}
  \item \feature{tsfresh__h__maximum}
  \item \feature{tsfresh__v__standard_deviation}
  \item \feature{tsfresh__v__root_mean_square}
\end{itemize}

这些项可以进入 PHM2012 的 top-10，但物理含义主要重复 RMS、standard deviation、variance、maximum 等手工统计量。因此 Step K 的结论不是“tsfresh 无效”，而是“在当前配置和当前任务上，tsfresh 没有提供新的主线特征族”。

XJTU-SY full-size \config{manual_tsfresh_basic} 被阻塞。当前 backend 会构造巨大的 long-format 表，Step G 估算约 604M 行，并在特征提取完成前被系统终止。因此本阶段不采用 tsfresh 作为 XJTU-SY 主线特征集，也不把失败 run 当作算法效果结论。

\subsection{特征集最终决策}

最终主线为：

\begin{center}
\config{manual_basic}
\end{center}

理由如下：

\begin{itemize}[leftmargin=2em]
  \item \config{manual_basic} 在 XJTU-SY 和 PHM2012 上都能完成三任务分析。
  \item 它覆盖 RUL、Health State、Early Fault 三类退化问题中最稳定的幅值/能量/离散度信号。
  \item 它避免了 XJTU-SY full-size tsfresh 的工程瓶颈。
  \item PHM2012 tsfresh 成功运行，但 top 特征多为手工统计量的重复表达。
\end{itemize}

延后或不采用：

\begin{itemize}[leftmargin=2em]
  \item \config{manual_tsfresh_basic}：XJTU-SY blocked；PHM2012 successful but redundant。
  \item \config{tsfresh_efficient}：本轮未评估，需等待明确的 memory-aware tsfresh 工程阶段。
\end{itemize}
